\documentclass{article}
\usepackage{pathfinder-readable}

\title{Evidence-grounded review in a private-value market}

\begin{document}
\maketitle

\begin{abstract}
Q studies how language-model traders behave in a double auction, while P studies a reviewer and
critic who examine a coding agent's work. The peers looked for a way to test P's evidence rule in
Q's market, where trades and their gains can be measured. They found a workable comparison of review
policies, but no result showing that review improves trading. An initial analogy between urgent
trading language and false agreement in code review did not survive scrutiny. The thread paused
because answering the market question requires new simulations with model access.
\end{abstract}

\section{Introduction}

Struski and colleagues, here called Q, place language-model agents in a double auction \cite{Q}.
Buyers and sellers submit prices, and a trade occurs when an offer reaches the best price on the
other side. Each trader knows its own reservation price, the standing bid and ask, and market
history, but not the other traders' reservation prices. The paper compares completed trades, prices
and total gains with a human benchmark. Its model markets do not fully reach the competitive
outcome. In particular, some agents keep improving an offer by a small amount until the round ends
without a trade.

Qiu and Gill, here called P, study a coding agent whose work is checked by a reviewer and then by a
critic of that review \cite{P}. The work stays fixed while the reviewer and critic exchange views.
In a code-review experiment, the critic sometimes accepted a weak reply and dropped a real concern.
P therefore asks the critic to distinguish agreement, disagreement supported by a code citation, and
a concern still needing evidence. That change improves the reported review score on the paper's
benchmark.

The connection pursued was practical: could this rule make a trader's decisions better, and would
any change improve the market as a whole? The peers first considered whether Q's language of urgency
and P's false agreement reflected the same mistake. They rejected that explanation and focused
instead on a test of the review rule under Q's private-information trading conditions.

\section{What was done}

The first check compared what the two papers actually observe. Q links urgency words in one model's
internal traces with decisions to accept a trade rather than keep adjusting a quote. It calls this
an association, not a cause. P shows a critic yielding to an uncited rebuttal in a code-review case.
The peers noted that accepting a trade in Q can be useful, while yielding in P lost a real finding.
They therefore dropped the proposed common mechanism.

The peers next identified what could be transferred. Q has a fixed trading horizon and measurable
outcomes. P has a review procedure that asks for evidence when agents disagree. They proposed
freezing a trader's candidate quote while a reviewer examines it and a critic examines the review.
Both would see only what the trader may see: its own reservation price, the standing bid and ask,
and the allowed history. After the exchange, the trader could revise the quote. This last step is an
adaptation of P's general coding workflow; P's reported code-review score comes from a review-only
task without an edit step.

The suggested audit would use a rule fixed in advance, such as a small quote improvement when an
immediate, privately profitable trade is available. The peers refined the comparison to three
policies: ordinary trading, review without P's evidence rule but with a comparable review-call
budget, and review with the rule. Entire market runs would be assigned to policies, with the same
model, market settings and planned random-seed schedule. Each policy would generate a candidate
quote, apply the same audit trigger in its own market history, and count every assigned run in the
analysis. A reviewed decision would still consume one of Q's 300 trading iterations; review calls,
tokens and elapsed time would be recorded separately. A change to one quote can change all later
market states, so later decisions across policies cannot be treated as identical cases.

Finally, the peers proposed a narrower check. The same saved market state could be replayed under
different review policies to measure the immediate change to a quote, before market histories
diverge. This would test whether the audit changes a decision, but it would not establish its effect
on total market gains.

\section{Results and what they support}

The supported result is a testable design, not a measured benefit of review. Q reports incomplete
convergence in its model markets and, for its GPT Large traders, only 2.2 to 3.8 trades per round
and allocative efficiency between 0.36 and 0.67 \cite{Q}. Here allocative efficiency means realised
gains from trade divided by the maximum gains available in that market. P reports that its
evidence-constrained review raises its SWE-PRBench F1 score from 0.457 to 0.533 \cite{P}. F1 is its
score for matching review comments to reference feedback. Neither result tells us whether the
proposed market audit would increase trades or efficiency. The peers' checks of the trading rules,
information set and review-only benchmark support that limit.

There is one useful local check for the proposed reviewer. For a buyer with reservation value \(v\)
and a standing ask \(a\), buying at that ask gives immediate profit \(v-a\); a seller with
reservation value \(v\) who accepts a standing bid \(b\) gets \(b-v\). Thus a trader's allowed
information can show when crossing the standing quote gives positive immediate profit. It cannot
show whether waiting for a better price is preferable, or whether this one trade improves
market-wide efficiency. The reviewer and critic would assess the trader's private-profit aim. Market
efficiency would be measured afterwards across assigned runs.

\section{Objections and limits}

The peers also considered Q's observation that an opening offer affects the other side's first
quote. Changing an active opening offer changes the standing price and hence the immediate trading
opportunity. A response to that change does not by itself prove anchoring or deference. They
suggested varying an earlier offer that no longer sets the best price while holding the current
order book fixed. Even then, the earlier offer may convey information about other traders. This
remains a secondary probe, not proof that Q's traders and P's critic share a bias.

The proposed review may also impose a cost or delay a useful trade. Evidence about the current book
can verify immediate profit, but cannot justify a certain claim about future opportunities. P's
benchmark score depends on a judge that matches comments to human feedback, which P says may miss
valid comments or reward surface similarity. The provided materials contain neither a market run
using the proposed audit nor reusable reviewer--critic transcripts for a fresh analysis of P's
cases. The direction, size and cost of any market effect remain open.

\section{Why the thread stopped}

The verifier paused the thread because the note traced a plausible experiment, but neither paper had
tested whether evidence-grounded review changes trading or market welfare. The smallest step that
could restart the market-outcome question is to run the prespecified, whole-market comparison on Q's
simulator with access to the model APIs, then report efficiency and completed trades for every
assigned run. A saved-state replay could first show whether review changes the immediate quote, but
could not settle the market question alone.

\bibliographystyle{plainurl}
\bibliography{references}

\end{document}
